\begin{textsample}{POS Dim 2 – al – Score 22.00 – al/al\_em\_3.txt}  \label{ex:f2_pos_255}
UM NOVO ENSINO MÉDIO A Base Nacional Comum Curricular para a Etapa do Ensino Médio foi aprovada pelo Conselho Nacional de Educação e \textbf{homologada} pelo Ministério da Educação em dezembro de 2018.
Esse documento define direitos e objetivos de aprendizagem no ensino médio e irá nortear a construção dos \textbf{currículos} dos estados, respeitada a diversidade regional, estadual e local.
A BNCC será a referência para formação de professores, avaliação, elaboração de conteúdos educacionais e critérios para a \textbf{oferta} de infraestrutura adequada para o pleno desenvolvimento da educação.
A Lei nº 9.394, de 20 de dezembro de 1996, que estabelece as \textbf{diretrizes} e bases da educação nacional, regulamentando a estrutura e o funcionamento da educação básica brasileira, aponta a necessidade de construção de uma Base Nacional Comum Curricular.
Com base nos marcos constitucionais, a LDB, no Inciso IV de seu \textbf{Artigo} 9º, preconiza que cabe à União : > estabelecer, em colaboração com os Estados, o Distrito Federal e \textbf{os} Municípios, competências e \textbf{diretrizes} para a Educação Infantil, o Ensino Fundamental e o Ensino Médio, que nortearão os \textbf{currículos} e seus conteúdos mínimos, de modo a assegurar formação básica comum. ( BRASIL, 1996 ) A relação entre o que deve ser comum e o que deve ser diverso fica clara nos \textbf{Artigos} 26 e 36 da LDB, respectivamente, quando determinam que : > os \textbf{currículos} da Educação Infantil, do Ensino Fundamental e do Ensino Médio devem ter base nacional comum, a ser complementada, em cada sistema de ensino e em cada estabelecimento escolar, por uma parte diversificada, exigida pelas características regionais e locais da sociedade, da cultura, da economia e dos educandos ( BRASIL, 1996 ; ênfase adicionada ). (... ) > “ O \textbf{currículo} do ensino médio será composto pela Base Nacional Comum Curricular e por \textbf{itinerários} formativos, que deverão ser organizados por meio da \textbf{oferta} de diferentes arranjos curriculares, conforme a relevância para o contexto local e a possibilidade dos sistemas de ensino (... ) Com a alteração feita pela Lei 13.415/2017, a LDB implementou mudanças para o Ensino Médio, como o aumento da \textbf{carga} \textbf{horária} mínima, a ampliação das escolas de tempo integral e a possibilidade de que todos os estudantes da etapa escolham caminhos de aprofundamento dos seus estudos.
Em 2018, a Resolução nº 3 CNE/CEB \textbf{atualiza} as \textbf{Diretrizes} Curriculares Nacionais para o Ensino Médio ( DCNEM ), normas que orientam e definem o planejamento dos \textbf{currículos} de escolas e sistemas de ensino.
Fundamentada nas \textbf{Diretrizes} Curriculares Nacionais para o Ensino Médio, a \textbf{Portaria} nº 1.432/2018 estabelece os Referenciais Curriculares para a Elaboração de \textbf{Itinerários} Formativos, com base nos 4 eixos estruturantes.
No Plano Nacional de Educação ( Lei 13.005/2014 ) são estabelecidas \textbf{diretrizes}, metas e estratégias para a \textbf{política} educacional dos próximos dez anos ( até 2024 ).
A renovação do Ensino Médio, com abordagens interdisciplinares e \textbf{currículos} escolares \textbf{flexíveis}, está entre os objetivos do PNE ( Meta 3 ).
Uma das estratégias para alcançar essa meta é estabelecer e implantar as \textbf{diretrizes} pedagógicas para a educação básica e a base nacional comum dos \textbf{currículos}, com direitos e objetivos de aprendizagem e desenvolvimento dos ( as ) alunos ( as ) para cada ano do ensino fundamental e médio.
A importância da base nacional comum curricular para a qualidade da Educação Básica em todas as etapas e modalidades é \textbf{reiterada} pelo Plano Estadual de Educação de Alagoas, Lei nº 7.795, de 22 de \textbf{janeiro} de 2016.
Todo esse arcabouço legal tem o objetivo de fundamentar a promoção das mudanças necessárias para que os estudantes do ensino médio sejam formados para a vida, sendo eles o centro da aprendizagem.
Para isso é fundamental que ocorra a atualização e flexibilização curricular, a partir do que propõe a Base Nacional Comum Curricular e os Referenciais para os \textbf{Itinerários} Formativos, com vistas a práticas pedagógicas mais interativas, inclusivas e diversificadas.
Com o objetivo de apoiar o processo de revisão ou elaboração e implementação dos \textbf{currículos} alinhados à BNCC foi \textbf{instituído}, pela \textbf{Portaria} nº 331/2018, o Programa de Apoio à Implementação da Base Nacional Comum Curricular – ProBNCC.
O Programa apóia as unidades federativas por meio de assistência financeira, via Plano de Ações Articuladas - PAR às Seduc, formação oferecida pelo MEC para \textbf{equipes} de \textbf{currículo} e \textbf{gestão} do Programa nos estados, e assistência \textbf{técnica}.
O desenvolvimento integral dos estudantes deve ter o foco da avaliação, pois, além dos conhecimentos adquiridos, é preciso levar em consideração também as suas habilidades, atitudes e valores.
No intuito de fomentar o desenvolvimento, a autonomia, a responsabilidade e o protagonismo dos estudantes, a avaliação deve ser um instrumento que possibilita que esses estudantes sejam agentes ativos de seu próprio processo de ensino-aprendizagem.

% matched lemmas: artigo, atualizar, carga, currículo, diretriz, equipe, flexível, gestão, homologar, horário, instituir, itinerário, janeiro, oferta, os, política, portaria, reiterar, técnico
\end{textsample}
